\documentclass[11pt]{article}
\usepackage[T1]{fontenc}
\usepackage{lmodern,amsmath,amssymb,amsthm,mathtools}
\usepackage[margin=1in]{geometry}
\usepackage{enumitem,microtype,needspace}
\usepackage[hidelinks]{hyperref}
\newtheorem{theorem}{Theorem}[section]
\newtheorem{lemma}[theorem]{Lemma}
\newtheorem{corollary}[theorem]{Corollary}
\newtheorem{proposition}[theorem]{Proposition}
\theoremstyle{definition}\newtheorem{definition}[theorem]{Definition}
\theoremstyle{remark}\newtheorem{remark}[theorem]{Remark}
\DeclareMathOperator{\rk}{rank}
\DeclareMathOperator{\cl}{cl}
\DeclareMathOperator{\Core}{Core}
\newcommand{\OPT}{\mathrm{OPT}}
\newcommand{\setminusf}{\setminus}
\title{Robust Dedicated Actuator Placement\newline via Matroid Spanning}
\author{Qichao Wang\\
Hebei University of Technology\\[0.6em]
Daoyu Dong\\
University of Electronic Science and Technology of China}
\hypersetup{pdftitle={Robust Dedicated Actuator Placement via Matroid Spanning},pdfauthor={Qichao Wang and Daoyu Dong}}
\date{November 6, 2026}
\begin{document}
\maketitle
\begin{abstract}
We study dedicated actuators installed at distinct states of a structured
linear system, where failures remove actuators but leave all states and state
edges intact. For unit costs and one failure, an exact source-quota
relaxation followed by maximum matching in a common-repair graph gives a
$3/2$-approximation on arbitrary directed state graphs in $O(n^3)$ graph
operations. The proof applies the Tutte--Berge formula to disjoint unions of
fundamental cocircuits. For multiple failures, we prove exact fixed-parameter
tractability in the maximum-matching deficiency plus the failure budget,
independently of the number of source strongly connected components. The
underlying theorem optimizes robust matroid spanning with arbitrarily many
disjoint unit-cost quotas. We also prove a weighted safe/vulnerable spanning
theorem in the independence-oracle model, with a control application on
strongly connected graphs. 
\end{abstract}

\section{Introduction}
A structurally controllable system can lose controllability when an actuator
fails. Redundant placement must protect both the matching condition and
accessibility from the installed inputs. These requirements have different
combinatorial forms: robust spanning in a matroid and lower quotas on source
strongly connected components. Keeping both requirements explicit permits
approximation and exact parameterized algorithms within the same model.

Our first result is a cubic-time $3/2$-approximation for one actuator failure
and unit costs (Theorem~\ref{thm:three-halves}). Its repair argument is valid
for arbitrary matroids. The exact graph relaxation, not a minimum ordinary
controllability placement, provides the necessary lower bound. A three-state
example attains the ratio for a permissible relaxation choice.

Our second result concerns arbitrary independence-oracle matroids with
nonnegative purchase weights and safe or vulnerable elements
(Theorem~\ref{thm:weighted}). Building on the uniform-set core method of
Bentert et al.~\cite{BFGL}, a weight-ordered recursion preserves an optimum
robust purchased set. The recursion continues into lower-rank closures even
when a large uniform set has been found. A classical set-pairs inequality
sharpens the deficient-hyperplane count. The parameter is rank plus failure
budget, and the objective charges every purchased element.

The third result is exact unit-cost optimization with arbitrarily many
disjoint quota groups (Theorem~\ref{thm:quota-fpt}). A bounded robust core
need not contain all elements needed for accessibility. We instead preserve
its quota-completion cost by capping the number of representatives of each
group. The resulting control algorithm is FPT in matching deficiency plus
failure budget on arbitrary digraphs
(Corollary~\ref{cor:general-control-fpt}); the number of source components is
not a parameter.

Po\c{c}as, Pequito and Ramos~\cite{PPR} study the robust minimal structural
controllability problem, establishing NP-hardness and reporting approximation
factors $5/2$ for general digraphs and $2$ for strongly connected digraphs.
We specify the actuator-failure conventions in Section~2 and discuss
related robust-assignment and fault-tolerant matroid algorithms in
Section~\ref{sec:scope}.

\section{Model and the dual matching matroid}
Let $G=(E,\mathcal A)$ be a nonempty finite directed state graph with $n=|E|$;
self-loops are allowed. Each prescribed nonzero entry of the state and input
matrices is a separate free parameter. Structural controllability means that
a controllable numerical realization exists, equivalently that almost every
realization is controllable. A placement $P\subseteq E$ installs one
dedicated input at each selected state, at most one per state. An actuator
failure removes the installed input while the state vertices and arcs
remain intact; the placement is fixed before failures occur.

Write $S_1,\ldots,S_s$ for the source strongly connected components (SCCs),
those having no entering state edge from another SCC. Form the state
bipartite graph with a left and right copy of $E$, and an edge $u_Lv_R$
for each arc $u\to v$. Let $T$ be the transversal matroid on the right
copies: its independent sets are precisely the state sets matchable using
state edges. Put $M=T^*$ and $d=n-\rk_T(E)$, the maximum-matching deficiency.
For $X\subseteq E$, duality gives
\begin{equation}\label{eq:dual-rank}
 \rk_M(X)=|X|+\rk_T(E\setminus X)-\rk_T(E).
\end{equation}
In particular, $X$ spans $M$ if and only if $E\setminus X$ is independent
in $T$. A rank oracle for $M$ is computable by bipartite matching.

The standard structural-controllability criterion combines accessibility
with a matching saturating all state vertices~\cite{Lin,PKA}. With dedicated
inputs it yields the following useful formulation.
\begin{proposition}\label{prop:control}
A placement $P$ tolerates every set of at most $k$ actuator failures if and
only if
\begin{equation}\label{eq:control}
 \rk_M(P\setminus F)=d\quad(F\subseteq P,\ |F|\le k),
 \qquad |P\cap S_i|\ge k+1\quad(1\le i\le s).
\end{equation}
\end{proposition}
\begin{proof}
For a fixed surviving placement $X$, matching every state in $E\setminus X$
by state edges and every state in $X$ by its own input gives a matching
saturating all states. Conversely, an input-augmented saturating matching
matches every state in $E\setminus X$ using state edges. This proves the
matching equivalence, and~\eqref{eq:dual-rank} gives spanning in $M$.
Accessibility holds exactly when every source SCC contains an input: every
component is reachable from a source component in the condensation DAG,
and no input outside a source component can reach it. The source SCCs
retain an input after every at most $k$ failures exactly when they initially
contain at least $k+1$. Applying both criteria to every survivor proves the
claim.
\end{proof}

We use standard finite-matroid notation. The closure of $X$ is
$\cl(X)=\{e:\rk(X\cup\{e\})=\rk(X)\}$. A hyperplane is a flat of
rank one below the matroid rank. A loop has rank zero; a coloop belongs to
every basis. The coloops of a restriction $M|P$ are precisely the elements
whose deletion lowers $\rk_M(P)$. A circuit is a minimal dependent set.
For a basis $B$ and $e\notin B$, the unique circuit contained in $B\cup\{e\}$
is its fundamental circuit, including the one-element circuit when $e$ is a
loop. We use these elementary matroid properties~\cite{Oxley}.

\section{A cubic time three halves approximation}
Throughout this section $k=1$ and every actuator has unit cost.

\subsection{An exact source quota relaxation}
Minimize $|P|$ subject to $P$ spanning $M$ and $|P\cap S_i|\ge2$ for
all $i$. Denote its optimum by $L$. Every robust placement is feasible for
this relaxation, so $L\le\OPT$.

It is solvable by an integral maximum flow. The complement $I=E\setminus P$
must be matchable using state edges and must satisfy
$|I\cap S_i|\le |S_i|-2$. Use a source, the two state copies, one group
node per source SCC, and a sink. Source-to-left-state edges have capacity
one; left-to-right state edges have capacity one. Each right state has a
capacity-one outgoing edge, to its source group if it has one, and directly
to the sink otherwise. The edge from group $i$ to the sink has capacity
$|S_i|-2$. An integral maximum flow gives a maximum-cardinality admissible
$I$, and therefore a minimum-cardinality $P$.

A singleton source SCC makes robust placement impossible. Conversely, if
all source SCCs have at least two vertices, every vertex has an incoming
state edge: this is immediate within a nonsingleton SCC, and a singleton
nonsource SCC has an entering edge. Thus $T$ has no loops and $M$ has no
coloops. The full placement $E$ spans after every single deletion and
satisfies all quotas. This proves that singleton source SCCs characterize
infeasibility for this one-failure model.

\subsection{Cocircuit unions and a global lower bound}
The next two lemmas concern an arbitrary rank-$d$ matroid $M$ with a
spanning set $P$. Let $D$ be the coloops of $M|P$. For $b\in D$ define
\[
 C_b=E\setminus\cl(P\setminus\{b\}).
\]
The set $C_b$ consists of the elements that restore spanning after deleting
$b$. It meets $P$ exactly in $\{b\}$.

\begin{lemma}\label{lem:union-rank}
For every $R\subseteq D$,
\[
 \rk\Bigl(E\setminus\bigcup_{b\in R}C_b\Bigr)=d-|R|.
\]
\end{lemma}
\begin{proof}
Choose a basis $B\subseteq P$. Every $b\in D$ belongs to $B$, and
$B\setminus\{b\}\subseteq P\setminus\{b\}$ are both of rank $d-1$.
Their closures are consequently equal. For nonempty $R$ we claim
\[
 \bigcap_{b\in R}\cl(B\setminus\{b\})=\cl(B\setminus R).
\]
For $e\notin B$, membership in $\cl(B\setminus\{b\})$ means that its
fundamental circuit omits $b$. Membership for every $b\in R$ therefore
means that this circuit is contained in $(B\setminus R)\cup\{e\}$,
which is precisely membership in $\cl(B\setminus R)$. Basis elements
satisfy the equality directly. This also handles loops. Taking ranks
proves the identity. For $R=\varnothing$ it is simply $\rk(E)=d$.
\end{proof}

Form a simple undirected graph $K$ on $D$. Put an edge $bc$ when some
$e\in E\setminus P$ belongs to both $C_b$ and $C_c$, and retain one such
witness for each edge. This graph is not necessarily bipartite. Let $\nu$
be its maximum matching cardinality.

\Needspace{6\baselineskip}
\begin{lemma}\label{lem:packing}
Every set $W$ spanning $M$ after any one deletion satisfies
$|W|\ge2(|D|-\nu)$.
\end{lemma}
\begin{proof}
Fix $U\subseteq D$ and let $R_1,\ldots,R_t$ be the nonempty connected
components of $K-U$. Put $Q_i=\bigcup_{b\in R_i}C_b$. These sets are
pairwise disjoint. An element of $P$ belongs to at most one $C_b$. A common
external element of $Q_i,Q_j$ would witness an edge between their components.
Lemma~\ref{lem:union-rank} gives $\rk(E\setminus Q_i)=d-|R_i|$.
Since adding an element raises rank by at most one, $W$ must contain at
least $|R_i|$ elements of $Q_i$ to span. Equality is impossible: deleting
one of these elements would leave rank at most
$(d-|R_i|)+(|R_i|-1)<d$. Thus
\[
 |W|\ge\sum_i(|R_i|+1)=|D|-|U|+t
 \ge |D|-|U|+\operatorname{odd}(K-U),
\]
where $\operatorname{odd}$ counts odd-order components. The classical
Tutte--Berge formula~\cite{LP} states
\[
 2\nu=|D|-\max_{U\subseteq D}
 \bigl(\operatorname{odd}(K-U)-|U|\bigr).
\]
Taking the maximum in the preceding inequality proves the result. The
argument compares against every robust $W$, not just supersets of $P$.
\end{proof}

\subsection{Algorithm and guarantee}
\begin{theorem}\label{thm:three-halves}
For any nonempty directed state graph in the stated dedicated-actuator
model, a robust placement of size at most $3\OPT/2$ can be found, or
infeasibility reported, in $O(n^3)$ graph operations and $O(n^2)$ space.
The guarantee is for unit costs and one actuator failure.
\end{theorem}
\begin{proof}
Reject singleton source SCCs; otherwise compute an optimum relaxation $P$.
Find $D$ and a maximum matching of $K$. For each matched pair install its
external common-repair witness. For every unmatched $b\in D$, install any
external element of $C_b$. Such an element exists because $M$ has no
coloops, while $P\setminus\{b\}$ has rank $d-1$. Let $R$ be the set of
installed witnesses, with repetitions removed. There are at most
$\nu+(|D|-2\nu)=|D|-\nu$ witnesses.

The set $P\cup R$ is robust. If a witness fails, $P$ still spans. If an
original noncritical element fails, the surviving part of $P$ spans. If
$b\in D$ fails, its witness repairs the deficiency. The source quotas
persist under additions. By Lemma~\ref{lem:packing},
\[
 |P\cup R|\le L+|D|-\nu\le\OPT+\frac{\OPT}{2}.
\]
Witness reuse is legitimate and can only decrease the cardinality. When
$D=\varnothing$, $P$ itself is robust.

For the complexity, retain the state matching $J$ obtained from the quota
flow. It covers $I=E\setminus P$. For each $b\in P$, explore alternating
paths from unmatched right vertex $b$: traverse state edges to the left,
then matched edges back to the right. A reachable free left vertex is
exactly an augmenting path showing that $I\cup\{b\}$ is matchable, or
that $b$ is noncritical. For a critical $b$, the reachable matched right
vertices are exactly its external repairs. A path to right vertex $e$
can be flipped to match $I\setminus\{e\}\cup\{b\}$. Conversely, the
symmetric difference with any matching of that replacement set gives such
a path, since a path to a free left vertex would contradict criticality.
Hence these searches recover the complete repair supports.

There are $O(n)$ vertices and $O(n^2)$ edges in the flow network. At most
$n$ unit augmentations, each found in $O(n^2)$ operations, suffice. At most
$n$ alternating searches also cost $O(n^3)$. For each external state,
connect all pairs of critical states it repairs, storing a witness for each
edge. This clique construction costs $O(n^3)$. A maximum general-graph
matching is computable in $O(n^3)$ operations by a blossom algorithm
\cite{LP}; a maximal matching is insufficient. All structures require $O(n^2)$ space.
\end{proof}

For the bidirected three-state star with center $a$ and leaves $b,c$, the
relaxation can choose $P=\{a,c\}$. Its one critical state is $c$, and the
only external repair is $b$. The algorithm returns all three states, while
$\{b,c\}$ is robust. Thus the ratio $3/2$ is attained for this relaxation
choice. Choosing the optimum robust pair at the relaxation stage avoids this
worst case.

\section{Weighted robust matroid spanning}
We now leave the unit-cost control model and state an abstract theorem.
Let $M=(E,\mathcal I)$ be a rank-$r$ matroid given by an independence
oracle, let $V\subseteq E$ be its vulnerable elements, and let $k\ge0$.
A set $B$ is \emph{resilient} if
\begin{equation}\label{eq:resilience}
 \rk_M(B\setminus F)=r
 \quad\text{for every }F\subseteq B\cap V,\quad |F|\le k.
\end{equation}
Elements in $E\setminus V$ are safe. With nonnegative binary-encoded
weights $w$, minimize $w(B)=\sum_{e\in B}w(e)$. Failures concern the
selected set, not survival of the whole ambient matroid.

\begin{theorem}\label{thm:weighted}\label{thm:main}
An optimum resilient purchased set can be computed, or infeasibility
reported, in
$f(r,k)\operatorname{poly}(n,\ell)$ time and oracle calls, where $\ell$ is the
input bit length and independence queries have unit cost. For $r\ge2$,
writing $q=(k+1)r$, a bound is
\[
 2^{O(qr^2(r+k))}\operatorname{poly}(n,\ell),
\]
hence $2^{O((r+k)^5)}\operatorname{poly}(n,\ell)$. The cases $r\le1$ or
$k=0$ are polynomial-time solvable.
\end{theorem}
\subsection{Matroid preliminaries and a small optimum}
Loops can be discarded because their removal preserves resilience and
cannot increase nonnegative purchase cost. If $r=0$, the empty set is
optimal. If $k=0$, a minimum-weight basis is optimal by matroid greedy.
For $r=1$, choose the cheaper feasible option between a cheapest safe
nonloop and the $k+1$ cheapest vulnerable nonloops. Henceforth assume
$r\geq2$ and $k\geq1$, and work on the loop-free restriction. Set
\begin{equation}\label{eq:parameters}
 q=(k+1)r,\qquad H=\binom{r+k-1}{k},\qquad t_h=(h-1)H+q\quad(1\leq h\leq r).
\end{equation}
The same global values of $q,H$ are used in every recursive subproblem.

\begin{lemma}\label{lem:small}
Every resilient set $B$ contains a resilient subset $B_0$ of cardinality at most $q$. Consequently, if a feasible solution exists, there is a minimum-weight solution of cardinality at most $q$.
\end{lemma}
\begin{proof}
Successively choose a maximal independent subset $I_j$ of $B\setminus\bigcup_{i<j}I_i$, for $j=0,\ldots,k$, and put $B_0=\bigcup_{j=0}^{k}I_j$. Empty layers are allowed. Since $|I_j|\leq r$, we have $|B_0|\leq q$.

Fix $F\subseteq B_0\cap V$ with $|F|\leq k$. The layers are disjoint, so some $I_j$ avoids $F$. Every element of $B\setminus B_0$ was available when $I_j$ was chosen and belongs to $\cl(I_j)$. Every element of $B_0\setminus F$ already lies in that surviving set. Therefore
\[
 B\setminus F\subseteq \cl(B_0\setminus F).
\]
Since $F$ is also an allowed failure of $B$, its left-hand side spans $M$. Thus $B_0$ is resilient. Nonnegative weights give $w(B_0)\leq w(B)$, and applying this construction to an optimum proves the second assertion.
\end{proof}

\subsection{A weight respecting exchange}
We first sharpen the deficient-hyperplane count using the classical set-pairs inequality of Bollob\'as~\cite{Bollobas}. We use the following bound.

\begin{lemma}\label{lem:hyperplanes}
For any $Y\subseteq E$, at most $\binom{r+k-1}{k}$ distinct hyperplanes can have the form $\cl(Y\setminus F)$ with $F\subseteq Y$ and $|F|\leq k$. This count is sharp.
\end{lemma}
\begin{proof}
Index these hyperplanes by $P_i$. Choose a basis $A_i\subseteq Y\cap P_i$ of size $r-1$, and put $D_i=Y\setminus P_i$. Because $P_i=\cl(Y\setminus F)$ for some permitted $F$, we have $D_i\subseteq F$, so $|D_i|\leq k$. Also $A_i\cap D_i=\varnothing$. If $i\ne j$, then $A_i\cap D_j\ne\varnothing$: otherwise $A_i\subseteq P_j$ and hence $P_i=\cl(A_i)\subseteq P_j$, forcing equality of the two hyperplanes.

In a uniformly random permutation of $Y$, let $\mathcal E_i$ be the event that every element of $A_i$ precedes every element of $D_i$. Its probability is $\binom{r-1+|D_i|}{r-1}^{-1}$. These events are pairwise disjoint. Indeed, for $i\ne j$ choose $a\in A_i\cap D_j$ and $b\in A_j\cap D_i$. The two events require $a$ before $b$ and $b$ before $a$, respectively. Therefore
\[
 \sum_i\binom{r-1+|D_i|}{r-1}^{-1}\leq1.
\]
Each summand is at least $\binom{r+k-1}{r-1}^{-1}$, proving the bound. For sharpness take $M=U_{r,r+k}$, $B=E$, all elements vulnerable, and $Y=B\setminus\{x\}$. The set $B$ is resilient, and each $k$-subset of $Y$ leaves a distinct $(r-1)$-element hyperplane. Thus equality holds even for the deficient hyperplanes needed below.
\end{proof}

Call $Z$ \emph{$h$-uniform} if $\rk(Z)=h$ and every $h$-element subset of $Z$ is independent. Such a set has at least $h$ elements, and every subset of size at most $h$ is independent. A set has a single failure status if it lies entirely in $V$ or entirely in $E\setminus V$.

\begin{lemma}\label{lem:exchange}
Let $A\subseteq E$ have a single failure status and rank $h\geq1$. Let $Z\subseteq A$ be $h$-uniform with $|Z|\geq t_h$. Suppose $B$ is resilient, $|B|\leq q$, and $x\in(B\cap A)\setminus Z$ satisfies $w(z)\leq w(x)$ for every $z\in Z$. Then some $y\in Z\setminus B$ makes
\[
 B'=(B\setminus\{x\})\cup\{y\}
\]
resilient, with $|B'|=|B|$ and $w(B')\leq w(B)$.
\end{lemma}
\begin{proof}
Put $Y=B\setminus\{x\}$. Consider the permitted failures $F\subseteq Y\cap V$, $|F|\leq k$, for which $\rk(Y\setminus F)<r$. Resilience of $B$ implies that their surviving ranks equal $r-1$ and that $x\notin\cl(Y\setminus F)$. Lemma~\ref{lem:hyperplanes} bounds the number of these distinct hyperplanes by $H$. It applies also when $x$ is safe and some deficient failure sets have size less than $k$.

Each of these hyperplanes contains at most $h-1$ elements of $Z$. Otherwise it contains an independent $h$-subset of $Z$, which spans $A$, and hence contains $x$, a contradiction. At most $q-1$ elements of $Z$ belong to $B$, because $x\notin Z$. The number of forbidden choices in $Z$ is at most $(h-1)H+q-1<t_h$. Choose a $y\in Z\setminus B$ outside all the deficient hyperplanes.

Let $F'\subseteq B'\cap V$ have size at most $k$. If $y\notin F'$, then $F'\subseteq Y$; either $Y\setminus F'$ already spans or $y$ repairs its rank deficiency by construction. If $y\in F'$, then $y$ is vulnerable, so $x$ is vulnerable as well. The set
\[
 (F'\setminus\{y\})\cup\{x\}
\]
is an allowed failure of $B$ of the same size as $F'$. Its survivor is exactly $B'\setminus F'$, which therefore spans. The cardinality and weight assertions follow from the choice of $y$.
\end{proof}

\subsection{Recursive construction of an optimum preserving core}
For a loop-free set $A$ of a single failure status, define $\Core(A)$ by the following procedure. Every scan uses a fixed total ordering that refines nondecreasing weight.
\begin{enumerate}[label=\arabic*.,leftmargin=*]
\item If $A=\varnothing$, return $\varnothing$. Set $h=\rk(A)$. If $|A|\leq t_h$, return $A$.
\item Scan the elements of $A$ in order, initially with $Z=\varnothing$. Accept an element $x$ precisely when every subset of $Z\cup\{x\}$ of size at most $h$ is independent. Stop when $|Z|=t_h$ or all elements have been scanned.
\item For every $T\in\binom{Z}{h-1}$, put $A_T=A\cap\cl(T)$ and recursively compute $\Core(A_T)$. Return
\[
 Z\ \cup\!\!\bigcup_{T\in\binom{Z}{h-1}}\Core(A_T).
\]
\end{enumerate}
In particular, step~3 is executed under \emph{both} stopping conditions. Recursive domains need not be flats of the original matroid.

Until $Z$ has size $h$, the acceptance rule is ordinary greedy independence. The final $Z$ therefore has at least $h$ elements and is $h$-uniform. Each child has rank exactly $h-1$: it contains $T$ and is contained in $\cl(T)$.

\begin{lemma}\label{lem:core}
For every resilient $B$ with $|B|\leq q$ and every single-status loop-free set $A$, there is a resilient $B'$ such that
\[
 |B'|=|B|,\quad w(B')\leq w(B),\quad B'\setminus A=B\setminus A,
 \quad B'\cap A\subseteq\Core(A).
\]
\end{lemma}
\begin{proof}
Induct on $h=\rk(A)$. Empty domains and a return of all of $A$ require no change. For $h=1$, every scanned nonloop is accepted, so the threshold scan
retains the $q$ lightest elements. Lemma~\ref{lem:exchange} replaces any
selected element outside $Z$ by a retained element of no greater weight;
the sole child domain is $A\cap\cl(\varnothing)=\varnothing$.
This establishes the base case.

Suppose $h>1$ and the scan is executed. Let $U=\bigcup_T A_T$. Every rejected scanned element $x$ belongs to $U$. Indeed, at rejection a dependent subset of $Z\cup\{x\}$ contains $x$ and at most $h-1$ retained elements. Those retained elements are independent, so they span $x$. Extend them to an $(h-1)$-subset of the final $Z$, obtaining a child containing $x$. Also $Z\subseteq U$, since $h>1$ and every element of $Z$ belongs to an $(h-1)$-subset of $Z$.

If the scan is exhausted, $A=U$. Otherwise each element of $A\setminus U$ is unscanned and is at least as expensive as every element of $Z$. Apply Lemma~\ref{lem:exchange} successively to the elements of $B\cap(A\setminus U)$. Each replacement lies in $Z\subseteq U$ and removes one selected element outside $U$. The resulting resilient set has its intersection with $A$ inside $U$, unchanged cardinality and outside-$A$ part, and no larger weight.

Fix the final output $W_A=\Core(A)$ and process the children $A_1,\ldots,A_m$ in any order. Apply the induction hypothesis to each child and the current solution. For the current solution define $D(B)=(B\cap A)\setminus W_A$. Processing $A_i$ changes only elements in $A_i$ and leaves its selected elements in $\Core(A_i)\subseteq W_A$. Consequently,
\[
 D(B_{\mathrm{new}})\subseteq D(B_{\mathrm{old}})\setminus A_i.
\]
Initially $D(B)\subseteq U$. After all children have been processed it is empty, regardless of overlaps between children. All operations preserve cardinality, hence the bound $q$, and preserve resilience without increasing weight. This establishes the induction.
\end{proof}

Apply the construction separately to the vulnerable and safe nonloop elements, and let $W$ be the union of the two outputs. Their domains are disjoint. Lemmas~\ref{lem:small} and~\ref{lem:core} imply that $W$ contains an optimum of cardinality at most $q$ whenever the instance is feasible.

\subsection{Size and running time}
Put $K_0=0$, $K_1=q$, and, for $2\leq h\leq r$, define
\[
 K_h=t_h+\binom{t_h}{h-1}K_{h-1}.
\]
By induction, a rank-$h$ call returns at most $K_h$ elements. The sequence is nondecreasing; in particular $|W|\leq2K_r$, or $|W|\leq K_r$ if all elements are vulnerable.

For a convenient explicit bound put $T=t_r$ and $N=2rT^{(r-1)^2}$. Each of the two recursion trees has at most $r$ levels and branching at most $T^{r-1}$. Thus there are at most $N$ nodes and $|W|\leq NT$. At a node, before $h$ elements are accepted, one independence query tests a candidate. Subsequently it suffices to test the $h$-subsets containing the candidate. This uses at most $O(n(1+T^{r-1}))$ independence queries per node. Testing $x\in\cl(T')$ for a child-defining independent $(h-1)$-set $T'$ uses one independence query to $T'\cup\{x\}$, with the membership case $x\in T'$ handled directly. Computing a domain's rank requires at most $n$ queries. Sorting and set bookkeeping have polynomial input cost multiplied by a parameter-only factor.

Finally enumerate all $B\subseteq W$ with $|B|\leq q$, and for each enumerate all allowed failure sets and compute its surviving ranks greedily. There are at most $(NT+1)^q$ candidates and at most $2^q$ failures per candidate. Choose a feasible candidate of minimum total weight. If there is none, optimum preservation proves infeasibility. Exact weight sums and comparisons have polynomial bit cost in the original weight encoding.

For $r\geq2$, put $p=r+k$. Since $H\leq2^{p-1}$ and $q\leq(p+1)^2$, we have $\log_2 T=O(p)$. Hence $NT=2^{O(r^2p)}$, and enumerating the candidates and failures gives $2^{O(qr^2p)}\operatorname{poly}(n,\ell)$. As $q=O(p^2)$, this is at most $2^{O(p^5)}\operatorname{poly}(n,\ell)$. For fixed $k\geq1$, the sharper estimate $H=O_k(r^k)$ gives $T=O_k(r^{k+1})$, $NT=r^{O_k(r^2)}$, and running time $r^{O_k(r^3)}\operatorname{poly}(n,\ell)$.

Nonnegative rational weights can also be handled by exact comparisons
and sums of polynomial bit complexity.

\subsection{Why simpler weighted reductions do not suffice}
\paragraph{Stopping at a large uniform set.}
Consider a rank-two truncation of a partition matroid with two parallel classes of size two having weights $(0,0)$ and $(1,1)$, and five further singleton classes each of weight $100$. Let $k=1$. Then $q=4$, $H=2$, and $t_2=6$. The weight-ordered uniform scan reaches the threshold while retaining only one element from each cheap parallel class. Its robust optimum is $101$, whereas the full instance has optimum $2$, attained by the two cheap pairs. Thus sorting the scan is insufficient if a successful threshold causes the recursion to stop. Our construction retains the cheap rejected elements through the lower-rank children.

\paragraph{Preserving cheapest ordinary bases after failures.}
Dey and Kavitha~\cite{DK} study a preserver $S\subseteq E$ such that, for every $F\subseteq E$ with $|F|\leq k$, $M|(S\setminus F)$ contains a minimum-cost ordinary basis of $M|(E\setminus F)$. That property need not preserve the optimum total cost of a resilient purchased set.

For a five-element example, take a rank-two truncation with parallel classes $A,B$ of size two and a singleton class $C$. Give each element of $A\cup B$ weight $2$ and the element of $C$ weight $3$. For every zero or one ground-element failure, $A\cup B$ still contains an ordinary basis of cost $4$, which is minimum in the full surviving matroid. Yet the robust optimum inside $A\cup B$ is $8$, and the full instance has optimum $7$, using one element from each class. Two successive cheapest residual bases also produce $A\cup B$. Thus this preserver construction does not replace the optimum-preserving recursion above.

\subsection{An application to dedicated input placement}
Consider a nonempty strongly connected state digraph with nonnegative
actuator costs and specified safe or vulnerable locations. Keep the
single-actuator-per-state and actuator-only failure conventions of Section~2.
Its dual state matching matroid has rank equal to the deficiency $d$;
Proposition~\ref{prop:control} supplies the matching interpretation, also
when only vulnerable actuators may fail.

\begin{corollary}
Minimum-cost robust dedicated-input placement on a strongly connected state digraph, under the actuator-failure model just specified, is FPT in $d+k$, with nonnegative binary-encoded costs and safe/vulnerable actuator locations.
\end{corollary}
\begin{proof}
If $d>0$, every surviving spanning set in $M$ is nonempty. Strong connectivity then ensures accessibility, so Theorem~\ref{thm:main} applies. If $d=0$, matching is automatic and accessibility requires a nonempty survivor. Choose the cheaper of one cheapest safe actuator and $k+1$ cheapest vulnerable actuators, omitting unavailable options.
\end{proof}
For general state digraphs, additional source SCC constraints are necessary. They are not preserved by the weighted same-status exchange alone. Section~\ref{sec:quota-fpt} supplies a separate unit-cost quota theorem, not a general weighted extension.


\Needspace{14\baselineskip}
\section{Exact parameterized optimization with source quotas}
\label{sec:quota-fpt}

The spanning requirement alone does not enforce accessibility in a general
state digraph. We now retain the source-SCC quotas without adding their number
to the parameter. Throughout this section every selected element is vulnerable
and costs one.

\begin{theorem}\label{thm:quota-fpt}
Let $M$ be a rank-$d$ matroid on $E$, given by an independence oracle, and let
$S_1,\ldots,S_s$ be pairwise disjoint subsets of $E$. For $k\ge0$, a
minimum-cardinality set $B$ satisfying
\[
 \rk_M(B\setminus F)=d\quad(F\subseteq B,\ |F|\le k),
 \qquad |B\cap S_i|\ge k+1\quad(1\le i\le s)
\]
can be found, or infeasibility certified, in $f(d,k)$ times a polynomial
in the input size and oracle cost. The parameter does not include $s$.
\end{theorem}

Set $\beta=k+1$, $G_0=E\setminus\bigcup_i S_i$, and $q=\beta d$. An undersized
group makes the instance infeasible. If $d=0$, choosing $\beta$ elements from
each group is optimal. Assume henceforth that $d\ge1$.

By Lemma~\ref{lem:small}, every robust spanning set contains a robust spanning subset of size at most
$q$: take $k+1$ successive maximal independent sets from its remaining
elements. Every omitted element is spanned by every layer, and one layer
avoids any $k$ failures. The union of the layers therefore remains spanning.
This small core contains no loops, although loops remain available for
quota completion. For a core $C$ define
\[
 \phi(C)=|C\cap G_0|+\sum_{i=1}^s\max\{0,|C\cap S_i|-\beta\}.
\]
Its cheapest quota completion has cardinality $\beta s+\phi(C)$ and remains
robust. Conversely, a small core inside an optimum placement has completion
cost no greater than that optimum. It suffices, therefore, to minimize
$\phi$ over robust spanning sets of size at most $q$.

The ordinary robust core alone need not preserve completion cost. In
$U_{1,4}$ on $\{a,b,c,d\}$, with $k=1$ and the sole group
$S_1=\{c,d\}$, the robust core $\{a,b\}$ completes to four elements,
whereas the robust optimum $\{c,d\}$ has size two. The following color
argument retains the necessary group choices.

\paragraph{Exchange candidates.}
Set $H=\binom{d+k-1}{k}$ as in Lemma~\ref{lem:hyperplanes}. Let $A$ be a loopless rank-$h$ subset and let
$Z\subseteq A$ be $h$-uniform: all of its subsets of size at most $h$ are
independent. Assume $|Z|\ge h$. For a robust set $B$ of size at most $q$
and $x\in(B\cap A)\setminus Z$, write $Y=B\setminus\{x\}$.
Whenever $\rk_M(Y\setminus F)<d$ for $F\subseteq Y$, $|F|\le k$,
the closure of $Y\setminus F$ is a hyperplane excluding $x$. There are
at most $H$ such hyperplanes, since each has a $(d-1)$-element basis in
$Y$. Each contains at most $h-1$ points of $Z$; otherwise it would contain
$A$ and hence $x$.
Thus at most $(h-1)H+q-1$ points of $Z$ lie either in a deficient
hyperplane or in $Y$. Every other $y$ gives a robust exchange
$Y\cup\{y\}$: failures avoiding $y$ are repaired by $y$, while a
failure containing $y$ corresponds to the failure of $x$ in $B$.

\paragraph{A color-capped core.} Start the recursion on the nonloop set $E^\circ=\{e\in E:\rk_M(\{e\})=1\}$. Keep all original elements, including loops, available for quota completion.
Use $G_0,S_1,\ldots,S_s$ as colors and set
\[
 c_h=(h-1)H+q,\qquad T_h=(q+1)c_h.
\]
In a rank-$h$ set $A$, greedily scan its elements into an $h$-uniform
set $Z$, allowing at most $c_h$ points of each color. Stop at $|Z|=T_h$
or at exhaustion. For $h=1$ return $Z$. Otherwise return $Z$ together
with recursively computed cores of
\[
 A\cap\operatorname{cl}_M(T),\qquad T\in\binom{Z}{h-1}.
\]
The recursion is performed in both stopping cases. The scan attains at
least $h$ points: before that the color cap cannot bind because
$c_h\ge q\ge d\ge h$, so the scan is ordinary greedy independence.
Every child consequently has rank exactly $h-1$.

We prove the following strong invariant by induction on $h$: every robust
$B$ of size at most $q$ can be replaced, without changing its size or its
part outside $A$, by a robust $B'$ whose $A$-part lies in the returned
core and for which $\phi(B')\le\phi(B)$.
Let $U=Z$ for $h=1$; otherwise let $U$ be the union of the child ground
sets $A\cap\operatorname{cl}_M(T)$. Then $Z\subseteq U$. A scanned
point rejected for dependence lies in $U$: its witnessing dependent set
shows that it is spanned by at most $h-1$ retained points, which can be
padded to one of the sets $T$.

Consider $x\in(B\cap A)\setminus U$. If it was scanned, it was rejected
by its color cap. The $c_h$ same-color points of $Z$ contain a valid
exchange candidate, since at most $c_h-1$ are forbidden by the preceding
count. The exchange preserves all color counts, and hence $\phi$.
If $x$ was unscanned, the total threshold was reached. With
$Y=B\setminus\{x\}$, forbid $G_0$ and all groups satisfying
$|Y\cap S_i|\ge\beta$. There are at most $q$ forbidden colors, since
$|Y|\le q-1$ and $b\ge1$. They contain at most $qc_h$ points of $Z$,
leaving at least $c_h$ candidates. A valid exchange therefore exists in
an unsaturated actual group. Its insertion has zero marginal $\phi$,
while deletion of $x$ cannot increase $\phi$.

These exchanges move the entire $A$-part into $U$. For $h=1$ this
finishes the proof. For higher rank, apply induction successively to
the children. If $W$ is the final returned core, a child operation cannot
create a point of $(B\cap A)\setminus W$ and removes every such point
in that child. All initially bad points lie in the union of the child
ground sets, so none remain after all children are processed. This also
justifies overlapping children and proves the invariant.

\paragraph{Size and running time.}
The output size is bounded by $L_d$, where
\[
 L_1=T_1,\qquad
 L_h=T_h+\binom{T_h}{h-1}L_{h-1}\quad(2\le h\le d).
\]
All these quantities depend only on $d,k$. Recursion depth is at most
$d$, branching is parameter bounded, and each node scans the input and
constructs its children using parameter bounded sets in independence
tests. Hence its oracle and bit cost is $f(d,k)\operatorname{poly}(|E|)$.
Enumerate subsets of the output of size at most $q$, check all at most
$2^q$ failures, minimize $\phi$, and complete the quotas using original
elements. The invariant preserves an optimum, and the enumeration has
parameter bounded size. This proves Theorem~\ref{thm:quota-fpt}.

\Needspace{8\baselineskip}
\begin{corollary}\label{cor:general-control-fpt}
For arbitrary directed state graphs, minimum unit-cost dedicated-actuator
placement tolerating at most $k$ actuator failures is fixed-parameter
tractable in the maximum-matching deficiency $d$ plus $k$, with at most
one actuator per state.
\end{corollary}
\begin{proof}
Use the dual of the state matching transversal matroid, whose rank is $d$.
Surviving matching feasibility is spanning in this dual, and surviving
accessibility is equivalent to selecting at least $k+1$ states in each
source SCC. The latter sets are disjoint, so Theorem~\ref{thm:quota-fpt}
applies.
\end{proof}

The returned core need not contain a complete feasible actuator placement:
quota completion may require many original elements. Thus the result is an
FPT algorithm, not a claim of a bounded full-solution kernel.

\section{Relationship to earlier work and limits}\label{sec:scope}
The robust purchased-set objective differs from choosing an ordinary basis
after failures. Bentert et al.~\cite{BFGL} prove unweighted FPT in rank plus
failure budget and explicitly pose the weighted question. Their proceedings
Proposition~6 gives the layered cardinality bound, Lemma~8 the uniform-set
exchange and basis-subset hyperplane count, and Lemma~9 the recursive core.
These are inherited ingredients. Our weighted argument adds the ordered
exchange invariant, recursion in both stopping cases, and separation of
failure statuses. The sharper count is an application of Bollob\'as's
classical set-pairs inequality~\cite{Bollobas}.
The same predecessor already gives polynomial weighted algorithms for rank
at most two and for truncations of unit-capacity partition matroids.

The all-vulnerable case of Theorem~\ref{thm:weighted} matches the weighted
problem defined in Section~1.1 of~\cite{BFGL}. Its Section~7 nonuniform
question prints $\rk(M-F)=\rk(M)$; the selected-set formulation is
$\rk_M(B\setminus F)=\rk(M)$. In personal correspondence with Qichao Wang
on 5 November 2026, Matthias Bentert confirmed that the latter was intended
and the former was a typographical error. This confirms the definition,
not the proofs in this paper. The two
printed predicates are not equivalent: in $U_{1,2}$, failure destroys the
selected singleton $B=\{a\}$ but does not lower the rank of the full
matroid. All our statements use the explicit selected-set definition
\eqref{eq:resilience}. This selected-object interpretation also appears in
flexible fault-tolerant network design~\cite{AHMS}.

Adjiashvili, Bindewald and Michaels~\cite{ABMnodes}, Theorem~7, give a
$7/4$-approximation for uniform cardinality node-robust assignment. Their
repair perspective is a predecessor of the present approach. Their
ICALP 2016 Theorem~5~\cite{ABMedges} already gives $3/2$ for uniform
\emph{edge}-robust assignment, buying assignment edges and tolerating an
edge failure via ear decompositions. That is a different purchase/failure
model. The proof in Section~3 instead uses fundamental cocircuit unions,
maximum common-repair matching, and source quotas.

Bulk-robust optimization~\cite{ASZ} charges purchased elements but supplies
an explicit family of failure scenarios. Expanding all at most $k$ failures
can require $n^{\Theta(k)}$ scenarios, so polynomial time in that expanded
list does not establish the present FPT bound. Weighted representative
families~\cite{FLPS} preserve independent extensions under their stated
representation and matroid hypotheses; their predicate is not directly
optimum preservation for a dependent purchased set under all failures.
Likewise, the cheapest-surviving-basis preservers of Dey and Kavitha~\cite{DK}
optimize a different object, as the example above shows.

The weighted theorem applies to safe and vulnerable elements, whereas the
quota theorem assumes unit costs and that all selected elements may fail.
Extending the latter to weighted quotas is an open direction. The optimal
polynomial-time approximation factor also remains open: the three-state
equality example concerns our relaxation choice, not an inapproximability
bound for the problem.

\paragraph{Reproducibility.}
The accompanying repository provides dependency-free implementations and
independent finite checks. The \href{https://github.com/dongxuelian2/robust-structural-controllability/tree/b510683764600ccd318eb7db566c35a816c2edc2}{approximation code}
and the \href{https://github.com/dongxuelian2/robust-structural-controllability/tree/9d62c26633643d05973e3ecefc5d36e7a8b426c1}{weighted and quota code}
are available at fixed commits; the accompanying guide gives the commands.
The code exposes rank oracles, which can be simulated from independence
queries by greedy basis construction.

\Needspace{5\baselineskip}
\paragraph{AI assistance.}
OpenAI tools assisted with developing and checking arguments, writing
verification code, and preparing this manuscript.

\begin{thebibliography}{99}
\small
\setlength{\itemsep}{3pt}
\raggedright
\bibitem{ABMnodes} D. Adjiashvili, V. Bindewald, and D. Michaels.
Robust Assignments with Vulnerable Nodes.
\href{https://arxiv.org/abs/1703.06074}{\mbox{arXiv:1703.06074}}, 2017.
\bibitem{ABMedges} D. Adjiashvili, V. Bindewald, and D. Michaels.
Robust Assignments via Ear Decompositions and Randomized Rounding.
In \emph{ICALP 2016}, LIPIcs 55, 71:1--71:14, 2016.
\href{https://doi.org/10.4230/LIPIcs.ICALP.2016.71}{doi:10.4230/LIPIcs.ICALP.2016.71}.
\bibitem{AHMS} D. Adjiashvili, F. Hommelsheim, M. M\"uhlenthaler, and O. Schaudt.
Fault-Tolerant Edge-Disjoint $s$-$t$ \mbox{Paths} -- Beyond Uniform Faults.
In \emph{SWAT 2022}, LIPIcs 227, 5:1--5:19, 2022.
\href{https://doi.org/10.4230/LIPIcs.SWAT.2022.5}{doi:10.4230/LIPIcs.SWAT.2022.5}.
\bibitem{ASZ} D. Adjiashvili, S. Stiller, and R. Zenklusen.
Bulk-Robust Combinatorial Optimization.
\emph{Mathematical Programming} 149, 361--390, 2015.
\href{https://doi.org/10.1007/s10107-014-0760-6}{doi:10.1007/s10107-014-0760-6}.
\bibitem{BFGL} M. Bentert, F. V. Fomin, P. A. Golovach, and L. Morelle.
Fault-Tolerant Matroid Bases. In \emph{ESA 2025}, LIPIcs 351,
83:1--83:14, 2025.
\href{https://doi.org/10.4230/LIPIcs.ESA.2025.83}{doi:10.4230/LIPIcs.ESA.2025.83}.
Full version: \href{https://arxiv.org/abs/2506.22010}{arXiv:2506.22010}.
\bibitem{Bollobas} B. Bollob\'as. On generalized graphs.
\emph{Acta Mathematica Academiae Scientiarum Hungaricae} 16,
447--452, 1965.
\href{https://doi.org/10.1007/BF01904851}{doi:10.1007/BF01904851}.
\bibitem{DK} D. Dey and T. Kavitha.
Low-Cost Arborescence Under Edge Faults.
\href{https://arxiv.org/abs/2605.13800}{arXiv:2605.13800}, 2026.
\bibitem{FLPS} F. V. Fomin, D. Lokshtanov, F. Panolan, and S. Saurabh.
Efficient Computation of Representative Families with Applications in
Parameterized and Exact Algorithms.
\emph{Journal of the ACM} 63(4), Article 29, 2016.
\href{https://doi.org/10.1145/2886094}{doi:10.1145/2886094}.
\bibitem{Lin} C.-T. Lin. Structural controllability.
\emph{IEEE Transactions on Automatic Control} 19(3), 201--208, 1974.
\bibitem{LP} L. Lov\'asz and M. D. Plummer.
\emph{Matching Theory}. North-Holland, 1986.
\bibitem{Oxley} J. Oxley. \emph{Matroid Theory}, second edition.
Oxford University Press, 2011.
\bibitem{PKA} S. Pequito, S. Kar, and A. P. Aguiar.
A Framework for Structural Input/Output and Control Configuration Selection
in Large-Scale Systems.
\emph{IEEE Transactions on Automatic Control} 61(2), 303--318, 2016.
\href{https://doi.org/10.1109/TAC.2015.2437525}{doi:10.1109/TAC.2015.2437525}.
\bibitem{PPR} D. Po\c{c}as, S. Pequito, and G. Ramos.
On the NP-hardness of the robust minimal structural controllability problem.
\emph{Automatica} 190, Article 113074, 2026.
\href{https://doi.org/10.1016/j.automatica.2026.113074}{doi:10.1016/j.automatica.2026.113074}.
\end{thebibliography}
\end{document}
